\documentclass[12pt]{article}
\usepackage[utf8]{inputenc}
\usepackage[T1]{fontenc}
\usepackage{hyperref}
\title{Takia Language Resources for K-12 Curriculum Development}
\author{Compiled by Claude Sonnet 5 (claude-sonnet-5) for Andr\'as Kornai}
\date{}
\begin{document}
\maketitle

\section{Introduction}
Takia (ISO 639-3: tbc; Glottocode: taki1248) is an Austronesian (Oceanic) language of Karkar Island and nearby areas of Madang Province, Papua New Guinea. The prompt gives about 40,000 speakers, but the search described that figure as undated and uncertain. Takia is well known in linguistics for contact-induced restructuring under the neighbouring Papuan language Waskia, which Malcolm Ross calls metatypy. The search found scholarly material but no school materials.

\section{Primary Educational Resources}
No primers, textbooks or school programmes for Takia were found in search results.

\section{Linguistic and Reference Materials}
\begin{itemize}
\item Wikipedia, ``Takia language''. \url{https://en.wikipedia.org/wiki/Takia_language}
\item Wikipedia, ``Metatypy'': Takia--Waskia contact. \url{https://en.wikipedia.org/wiki/Metatypy}
\item Wikipedia, ``Waskia language'': the neighbouring Papuan language. \url{https://en.wikipedia.org/wiki/Waskia_language}
\item Max Planck EVA, Takia numerals. \url{https://lingweb.eva.mpg.de/channumerals/Takia.htm}
\item ANU Research Portal: ``Loanwords in Takia, an Oceanic language of Papua New Guinea'' (Ross, 2009). \url{https://researchportalplus.anu.edu.au/en/publications/loanwords-in-takia-an-oceanic-language-of-papua-new-guinea/}
\item LINGUIST List 16.1811: a review page found in the search. I did not check what work it reviews. \url{https://linguistlist.org/issues/16/1811/}
\end{itemize}
The search suggests that Ross's chapter ``Takia'' in \emph{The Oceanic Languages} (Lynch, Ross and Crowley, 2002) is the main published description. I saw no link to it, and no separate Takia dictionary or grammar turned up.

\section{Religious and Community Resources}
None found in search results.

\section{Digital and Audio Resources}
No audio, video or corpus resources for Takia were found.

\section{Curriculum Development Context}
Takia is reportedly widely spoken and used as a first language by the whole community, which is a strong base for native-language teaching. The search turned up no orthography, readers or teaching materials, and no confirmation of the speaker count. Developing a curriculum would need new materials built with speakers, working from Ross's description and current community orthography practice, and checking with Madang Province education authorities.
\end{document}